% initial-draft-2 glossary (from draft-initial). Started from the Draft 1 definitions (unchanged in Draft 1/)
% and extended for the full-paper scope.

\newglossaryentry{amat}{
  name={AMAT (average memory access time)},
  description={A weighted average of memory-access latencies that accounts for cache hits, cache misses, and the penalties of accessing lower levels of the memory hierarchy}
}

\newglossaryentry{cache-coherence}{
  name={cache coherence},
  description={The mechanisms and rules used to maintain a consistent view of shared data when multiple caches can hold copies of the same memory location}
}

\newglossaryentry{cache-hierarchy}{
  name={cache hierarchy},
  description={A set of cache levels organised to exploit locality by balancing capacity, access latency, and implementation cost}
}

\newglossaryentry{ccd}{
  name={CCD (Core Complex Die)},
  description={AMD terminology for a compute die containing processor cores and associated cache resources in selected Zen-based processors}
}

\newglossaryentry{chiplet}{
  name={chiplet},
  description={A separately fabricated die designed to operate with one or more other dies as part of a packaged system}
}

\newglossaryentry{hybrid-bonding}{
  name={hybrid bonding},
  description={A fine-pitch die-stacking technique that forms direct mechanical and electrical connections between aligned die surfaces}
}

\newglossaryentry{infinity-fabric}{
  name={Infinity Fabric},
  description={AMD's interconnect architecture for communication among processor components and, in relevant products, among separate dies}
}

\newglossaryentry{iod}{
  name={IOD (I/O Die)},
  description={AMD terminology for a die that centralises input/output, memory-interface, and supporting interconnect functions in selected chiplet-based Zen processors}
}

\newglossaryentry{llc}{
  name={LLC (last-level cache)},
  description={The final cache level checked before a request proceeds to main memory in the cache hierarchy under discussion}
}

\newglossaryentry{memory-hierarchy}{
  name={memory hierarchy},
  description={An organisation of storage levels with different capacities, costs, and access times, arranged to exploit temporal and spatial locality}
}

\newglossaryentry{monolithic-die}{
  name={monolithic die},
  description={A single semiconductor die on which the principal processor resources under discussion are integrated rather than partitioned across multiple dies}
}

\newglossaryentry{tsv}{
  name={TSV (through-silicon via)},
  description={A vertical electrical connection passing through a silicon die to support communication in a stacked-die assembly}
}

\newglossaryentry{v-cache}{
  name={3D V-Cache},
  description={AMD's name for a design that uses three-dimensional die integration to add stacked cache capacity above a processor die}
}

\newglossaryentry{ccx}{
  name={CCX (Core Complex)},
  description={AMD terminology for a group of cores that share one L3 cache; in Zen~2 a CCD contains two four-core CCXs, and in Zen~3 a CCD contains one eight-core CCX}
}

\newglossaryentry{cpi}{
  name={CPI (cycles per instruction)},
  description={The average number of clock cycles taken per dynamic instruction; its reciprocal is IPC}
}

\newglossaryentry{defect-density}{
  name={defect density},
  description={The average number of yield-killing manufacturing defects per unit area of silicon, written $D_0$ in simple yield models}
}

\newglossaryentry{ifop}{
  name={IFOP (Infinity Fabric On-Package)},
  description={The on-package serial link used by AMD to connect compute dies to the I/O die over short organic-package routes}
}

\newglossaryentry{ipc}{
  name={IPC (instructions per cycle)},
  description={The average number of dynamic instructions completed per clock cycle; the reciprocal of CPI}
}

\newglossaryentry{micro-bump}{
  name={micro-bump},
  description={A small solder-based bump used to connect stacked or adjacent dies; its pitch limits the achievable interconnect density compared with hybrid bonding}
}

\newglossaryentry{reticle-limit}{
  name={reticle limit},
  description={The maximum die area that a lithography system can pattern in a single exposure field, which bounds the size of a monolithic die}
}

\newglossaryentry{sram}{
  name={SRAM (static random-access memory)},
  description={A memory built from bistable cells that retain data while powered without refresh; the usual technology for on-chip caches}
}

\newglossaryentry{yield}{
  name={yield},
  description={The fraction of manufactured dies that function correctly; larger dies generally yield worse for a given defect density}
}

\newglossaryentry{d2d-hop}{
  name={die-to-die hop},
  description={The extra latency of a transfer between two dies in a package, such as between a CCD and the I/O die. In the simulations of Part~V a hop is one CCD--IOD round trip, by default AMD's 9\,ns}
}

\newglossaryentry{l3d}{
  name={L3D (L3 die)},
  description={AMD's name for the stacked SRAM die of 3D V-Cache, called the cache die in this paper. A CCD carrying one is a V-Cache CCD; a CCD without one is a plain CCD. A product built from V-Cache CCDs is a V-Cache part}
}

\newglossaryentry{tlb}{
  name={TLB (translation lookaside buffer)},
  description={A small cache of recent virtual-to-physical address translations, consulted on every memory access}
}

\newglossaryentry{mlp}{
  name={memory-level parallelism},
  description={The overlap of several outstanding cache misses, which lets an out-of-order core wait roughly once for several misses}
}
